\chapter{Complexity — Adder}
%%% generated by the blueprint pipeline from TCSlib.Complexity.CircuitComplexity.Adder
\section{Overview}

Declarations of the Lean module \texttt{TCSlib.Complexity.CircuitComplexity.Adder}.

\section{Declarations}

\begin{definition}[Binary value of a Boolean list]
\lean{BoolCircuit.bitsVal}
\label{BoolCircuit.bitsVal}
\leanok
For a finite list of Boolean digits, define its binary value by interpreting the first
digit as the least significant digit and each subsequent digit as the next power of two.
\end{definition}

\begin{definition}[Binary prefix value]
\lean{BoolCircuit.prefixVal}
\label{BoolCircuit.prefixVal}
\leanok
For a Boolean sequence $f:\mathbb N\to\{0,1\}$, define $V_0=0$ and, for each $i\ge0$,
set $V_{i+1}=V_i+f(i)2^i$. Thus $V_i$ is the binary value of the first $i$ bits, with
$f(0)$ least significant.
\end{definition}

\begin{theorem}[Initial prefix value]
\lean{BoolCircuit.prefixVal_zero}
\label{BoolCircuit.prefixVal_zero}
\leanok
\uses{BoolCircuit.prefixVal}
\difficulty{1}
For every function $f:\mathbb{N}\to\{\mathrm{false},\mathrm{true}\}$, the value of its
prefix sum at index $0$ is $0$.
\end{theorem}

\begin{theorem}[Successor formula for a prefix value]
\lean{BoolCircuit.prefixVal_succ}
\label{BoolCircuit.prefixVal_succ}
\leanok
\uses{BoolCircuit.prefixVal}
\difficulty{1}
For every Boolean sequence $f:\mathbb{N}\to\{\mathrm{false},\mathrm{true}\}$ and natural
number $i$, the weighted prefix value below $i+1$ equals the weighted prefix value below
$i$ plus $(f(i))\,2^i$, where a true Boolean has value $1$ and a false Boolean has value
$0$.
\end{theorem}

\begin{theorem}[Shifted successor formula for a prefix value]
\lean{BoolCircuit.prefixVal_succ'}
\label{BoolCircuit.prefixVal_succ'}
\leanok
\uses{BoolCircuit.prefixVal, BoolCircuit.prefixVal_succ}
\difficulty{4}
For every Boolean sequence $f:\mathbb{N}\to\{\mathrm{false},\mathrm{true}\}$ and natural
number $i$, its weighted prefix value below $i+1$ equals $(f(0)) + 2$ times the weighted
prefix value below $i$ of the shifted sequence $j\mapsto f(j+1)$, where a true Boolean
has value $1$ and a false Boolean has value $0$.
\end{theorem}

\begin{theorem}[Bound on a binary prefix value]
\lean{BoolCircuit.prefixVal_lt}
\label{BoolCircuit.prefixVal_lt}
\leanok
\uses{BoolCircuit.prefixVal, BoolCircuit.prefixVal_succ}
\difficulty{4}
For every Boolean sequence $f:\mathbb{N}\to\{\mathrm{false},\mathrm{true}\}$ and every
natural number $i$, its weighted prefix value on positions below $i$ is strictly less
than $2^i$.
\end{theorem}

\begin{theorem}[Prefix values determined by initial bits]
\lean{BoolCircuit.prefixVal_congr}
\label{BoolCircuit.prefixVal_congr}
\leanok
\uses{BoolCircuit.prefixVal, BoolCircuit.prefixVal_succ}
\difficulty{4}
Let $f,g:\mathbb{N}\to\{\mathrm{false},\mathrm{true}\}$ be Boolean sequences and let $i$
be a natural number. If $f(j)=g(j)$ for every $j<i$, then their weighted prefix values
on positions below $i$ are equal.
\end{theorem}

\begin{theorem}[Injectivity of binary prefix values]
\lean{BoolCircuit.prefixVal_inj}
\label{BoolCircuit.prefixVal_inj}
\leanok
\uses{BoolCircuit.prefixVal, BoolCircuit.prefixVal_lt, BoolCircuit.prefixVal_succ}
\difficulty{5}
Let $f,g:\mathbb{N}\to\{\mathrm{false},\mathrm{true}\}$ be Boolean sequences and let $i$
be a natural number. If their weighted prefix values on positions below $i$ are equal,
then $f(j)=g(j)$ for every $j<i$.
\end{theorem}

\begin{theorem}[Boolean-list value as a prefix sum]
\lean{BoolCircuit.bitsVal_eq_prefixVal}
\label{BoolCircuit.bitsVal_eq_prefixVal}
\leanok
\uses{BoolCircuit.bitsVal, BoolCircuit.prefixVal, BoolCircuit.prefixVal_succ'}
\difficulty{4}
For every finite Boolean list $l$, its value obtained by interpreting its entries as
base-two digits equals $\sum_{j=0}^{|l|-1} b_j2^j$, where $b_j$ is the $j$th entry of
$l$, interpreted as $0$ or $1$.
\end{theorem}

\begin{theorem}[Two definitions of Boolean-list value agree]
\lean{BoolCircuit.bitsVal_eq_logProg_bitsVal}
\label{BoolCircuit.bitsVal_eq_logProg_bitsVal}
\leanok
\uses{BoolCircuit.bitsVal, Complexity.LogProg.bitsVal}
\difficulty{4}
For every finite Boolean list, its value obtained by interpreting its entries as
base-two digits equals the value defined recursively by assigning the first bit weight
$1$ and doubling the value of the remaining list.
\end{theorem}

\begin{definition}[Four-gate exclusive OR construction]
\lean{BoolCircuit.xorGates}
\label{BoolCircuit.xorGates}
\leanok
\uses{BoolCircuit.DAGGate}
Given natural-number input labels $x,y$ and a label $L$, define an ordered list of four
gates: an OR gate reading $x,y$; an AND gate reading $x,y$; a NOT gate reading $L+1$;
and an AND gate reading $L,L+2$.
\end{definition}

\begin{theorem}[Running the XOR gate sequence]
\lean{BoolCircuit.runWith_xorGates}
\label{BoolCircuit.runWith_xorGates}
\leanok
\uses{BoolCircuit.DAGGate.eval, BoolCircuit.runWith, BoolCircuit.runWith_cons, BoolCircuit.runWith_nil, BoolCircuit.xorGates}
\difficulty{4}
Let $vs$ be a list of Boolean values, and let $x,y,L\in\mathbb{N}$. Suppose $|vs|=L$,
$x<L$, and $y<L$. Running the four-gate sequence that computes the OR of the values at
$x,y$, their AND, the negation of the value at index $L+1$, and the AND of the values at
indices $L,L+2$, starting from $vs$, appends, in order, $vs_x\lor vs_y$, $vs_x\land
vs_y$, $\neg(vs_x\land vs_y)$, and $vs_x\mathbin{\mathrm{xor}}vs_y$, where an
out-of-range list index has value false.
\end{theorem}

\begin{definition}[Full-adder gate construction]
\lean{BoolCircuit.fullAdderGates}
\label{BoolCircuit.fullAdderGates}
\leanok
\uses{BoolCircuit.DAGGate, BoolCircuit.xorGates}
Given natural-number labels $a,b,c_{\mathrm{in}},L$, define a gate list by concatenating
the four-gate exclusive OR construction on $a,b$ at label $L$, the four-gate
construction on $L+3,c_{\mathrm{in}}$ at label $L+4$, and a final OR gate reading
$L+1,L+5$.
\end{definition}

\begin{theorem}[Outputs of a full-adder circuit]
\lean{BoolCircuit.runWith_fullAdderGates}
\label{BoolCircuit.runWith_fullAdderGates}
\leanok
\uses{BoolCircuit.DAGGate.eval, BoolCircuit.fullAdderGates, BoolCircuit.runWith, BoolCircuit.runWith_append, BoolCircuit.runWith_cons, BoolCircuit.runWith_nil, BoolCircuit.runWith_xorGates}
\difficulty{5}
Let $v$ be a Boolean list of length $L$, and let $a,b,c<L$ be indices into it.
Successively evaluate the full-adder gate list with input indices $a,b,c$ and base index
$L$, appending each gate's output to $v$. The resulting list is $v$ followed, in order,
by $v_a\lor v_b$, $v_a\land v_b$, $\neg(v_a\land v_b)$, $v_a\mathbin{\mathrm{xor}}v_b$,
$(v_a\mathbin{\mathrm{xor}}v_b)\lor v_c$, $(v_a\mathbin{\mathrm{xor}}v_b)\land v_c$,
$\neg((v_a\mathbin{\mathrm{xor}}v_b)\land v_c)$,
$(v_a\mathbin{\mathrm{xor}}v_b)\mathbin{\mathrm{xor}}v_c$, and $(v_a\land
v_b)\lor((v_a\mathbin{\mathrm{xor}}v_b)\land v_c)$, where $v_j$ denotes the Boolean
value at index $j$.
\end{theorem}

\begin{definition}[Ripple-carry bit block]
\lean{BoolCircuit.RippleCarry.bitBlock}
\label{BoolCircuit.RippleCarry.bitBlock}
\leanok
\uses{BoolCircuit.DAGGate, BoolCircuit.fullAdderGates, BoolCircuit.xorGates}
Given natural-number labels $a,b,c,c_{\mathrm{in}},\mathrm{acc},L$, define a gate list
by concatenating the full-adder construction on $a,b,c_{\mathrm{in}}$ at label $L$, the
four-gate exclusive OR construction on $L+7,c$ at label $L+9$, then a NOT gate reading
$L+12$ and an AND gate reading $\mathrm{acc},L+13$.
\end{definition}

\begin{theorem}[Evaluation of a ripple-carry bit block]
\lean{BoolCircuit.RippleCarry.runWith_bitBlock}
\label{BoolCircuit.RippleCarry.runWith_bitBlock}
\leanok
\uses{BoolCircuit.DAGGate.eval, BoolCircuit.RippleCarry.bitBlock, BoolCircuit.runWith, BoolCircuit.runWith_append, BoolCircuit.runWith_cons, BoolCircuit.runWith_fullAdderGates, BoolCircuit.runWith_nil, BoolCircuit.runWith_xorGates}
\difficulty{4}
Let $vs$ be a Boolean list and let $a,b,c,cin,acc,L$ be natural numbers. Suppose $vs$
has length $L$ and each of $a,b,c,cin,acc$ is less than $L$. Running the bit block with
these indices appends $15$ values to $vs$. The resulting value at index $L+8$ is
$(vs[a]\land vs[b])\lor((vs[a]\mathbin{\mathrm{xor}}vs[b])\land vs[cin])$, where an
out-of-range list lookup is false. The value at index $L+14$ is
$vs[acc]\land\neg(((vs[a]\mathbin{\mathrm{xor}}vs[b])\mathbin{\mathrm{xor}}vs[cin])\mathbin{\mathrm{xor}}vs[c])$,
with the same convention for out-of-range lookups.
\end{theorem}

\begin{definition}[Carry input vertex indices]
\lean{BoolCircuit.RippleCarry.cinIdx}
\label{BoolCircuit.RippleCarry.cinIdx}
\leanok
For each natural number $n$, define the carry input index at position $0$ to be $n$, and
at position $i+1$ to be $n+2+15i+8$.
\end{definition}

\begin{definition}[Accumulator vertex indices]
\lean{BoolCircuit.RippleCarry.accIdx}
\label{BoolCircuit.RippleCarry.accIdx}
\leanok
For each natural number $n$, define the accumulator index at position $0$ to be $n+1$,
and at position $i+1$ to be $n+2+15i+14$.
\end{definition}

\begin{definition}[Recursive list of ripple-carry adder gates]
\lean{BoolCircuit.RippleCarry.adderGatesUpTo}
\label{BoolCircuit.RippleCarry.adderGatesUpTo}
\leanok
\uses{BoolCircuit.DAGGate, BoolCircuit.RippleCarry.accIdx, BoolCircuit.RippleCarry.bitBlock, BoolCircuit.RippleCarry.cinIdx, BoolCircuit.constGate}
For natural numbers $n$ and $k$, define a list of gates $A_i$ for each natural number
$i$ by $A_0=(	ext{false constant gate},	ext{ true constant gate})$. For each $i$, let
$A_{i+1}$ be the concatenation of $A_i$ with the bit block whose three data inputs are
$i$, $k+i$, and $2k+i$, whose carry input is $n$ when $i=0$ and $n+2+15j+8$ when
$i=j+1$, whose accumulator input is $n+1$ when $i=0$ and $n+2+15j+14$ when $i=j+1$, and
whose base index is $n+2+15i$.
\end{definition}

\begin{theorem}[Length of the ripple-carry gate list]
\lean{BoolCircuit.RippleCarry.length_adderGatesUpTo}
\label{BoolCircuit.RippleCarry.length_adderGatesUpTo}
\leanok
\uses{BoolCircuit.RippleCarry.adderGatesUpTo, BoolCircuit.RippleCarry.bitBlock, BoolCircuit.fullAdderGates, BoolCircuit.xorGates}
\difficulty{4}
For all natural numbers $n,k,i$, the gate list consisting of the initial two constant
gates followed by $i$ bit blocks has length $2+15i$.
\end{theorem}

\begin{definition}[Ripple adder gate list]
\lean{BoolCircuit.adderGates}
\label{BoolCircuit.adderGates}
\leanok
\uses{BoolCircuit.DAGGate, BoolCircuit.RippleCarry.accIdx, BoolCircuit.RippleCarry.adderGatesUpTo, BoolCircuit.RippleCarry.cinIdx}
For each natural number $n$, define the gate list to be the list of gates up to position
$	ext{⌊}n/3	ext{⌋}$ with parameter $k=	ext{⌊}n/3	ext{⌋}$, followed by a NOT gate reading
the carry input index at that position and an AND gate reading the accumulator index
there and the vertex $n+2+15	ext{⌊}n/3	ext{⌋}$.
\end{definition}

\begin{theorem}[Carry-in index bound]
\lean{BoolCircuit.cinIdx_lt}
\label{BoolCircuit.cinIdx_lt}
\leanok
\uses{BoolCircuit.RippleCarry.cinIdx}
\difficulty{4}
For all natural numbers $n$ and $i$, the index defined recursively by $c_n(0)=n$ and
$c_n(i+1)=n+2+15i+8$ satisfies $c_n(i)<n+2+15i$.
\end{theorem}

\begin{theorem}[Accumulator index bound]
\lean{BoolCircuit.accIdx_lt}
\label{BoolCircuit.accIdx_lt}
\leanok
\uses{BoolCircuit.RippleCarry.accIdx}
\difficulty{4}
For all natural numbers $n$ and $i$, the accumulator index at step $i$ is less than
$n+2+15i$.
\end{theorem}

\begin{theorem}[Acyclicity of a bit block]
\lean{BoolCircuit.gatesAcyclic_bitBlock}
\label{BoolCircuit.gatesAcyclic_bitBlock}
\leanok
\uses{BoolCircuit.GatesAcyclic, BoolCircuit.RippleCarry.bitBlock, BoolCircuit.fullAdderGates, BoolCircuit.xorGates}
\difficulty{3}
Let $a,b,c,cin,acc,L$ be natural numbers, with each of $a,b,c,cin,acc$ strictly less
than $L$. The gate list formed by the specified full-adder gates, followed by the
specified XOR gates and two final gates, is acyclic starting from index bound $L$.
\end{theorem}

\begin{theorem}[Fan-in of a bit block]
\lean{BoolCircuit.faninTwo_bitBlock}
\label{BoolCircuit.faninTwo_bitBlock}
\leanok
\uses{BoolCircuit.DAGGate.FaninTwo, BoolCircuit.DAGGate.WellFormed, BoolCircuit.RippleCarry.bitBlock, BoolCircuit.fullAdderGates, BoolCircuit.xorGates}
\difficulty{3}
Let $a,b,c,\mathrm{cin},\mathrm{acc},L$ be natural numbers with $a\ne b$, $c<L$,
$\mathrm{cin}<L$, and $\mathrm{acc}<L$. Every gate in the corresponding bit block is
well formed and has at most two inputs.
\end{theorem}

\begin{theorem}[Acyclicity of partial adder gates]
\lean{BoolCircuit.gatesAcyclic_adderGatesUpTo}
\label{BoolCircuit.gatesAcyclic_adderGatesUpTo}
\leanok
\uses{BoolCircuit.GatesAcyclic, BoolCircuit.RippleCarry.adderGatesUpTo, BoolCircuit.RippleCarry.length_adderGatesUpTo, BoolCircuit.accIdx_lt, BoolCircuit.cinIdx_lt, BoolCircuit.gatesAcyclic_append, BoolCircuit.gatesAcyclic_bitBlock}
\difficulty{4}
Let $n$ and $k$ be natural numbers with $3k\le n$, and let $i$ be a natural number. The
list formed by starting with two constant gates and appending $i$ successive bit blocks
according to the ripple-carry construction is acyclic: each input index read by the gate
at position $j$ is less than $n+j$.
\end{theorem}

\begin{theorem}[Fan-in of partial adder circuits]
\lean{BoolCircuit.faninTwo_adderGatesUpTo}
\label{BoolCircuit.faninTwo_adderGatesUpTo}
\leanok
\uses{BoolCircuit.DAGGate.FaninTwo, BoolCircuit.RippleCarry.adderGatesUpTo, BoolCircuit.accIdx_lt, BoolCircuit.cinIdx_lt, BoolCircuit.faninTwo_bitBlock}
\difficulty{4}
Let $n,k,i$ be natural numbers with $3k\le n$ and $i\le k$. Every gate in the adder
circuit constructed through step $i$ is well formed and has at most two inputs.
\end{theorem}

\begin{theorem}[Acyclicity of the ripple-carry adder gates]
\lean{BoolCircuit.gatesAcyclic_adderGates}
\label{BoolCircuit.gatesAcyclic_adderGates}
\leanok
\uses{BoolCircuit.GatesAcyclic, BoolCircuit.accIdx_lt, BoolCircuit.adderGates, BoolCircuit.cinIdx_lt, BoolCircuit.gatesAcyclic_adderGatesUpTo, BoolCircuit.gatesAcyclic_append}
\difficulty{4}
For every natural number $n$, the list of gates obtained by starting with the two
constant gates, appending $n/3$ successive bit blocks with parameters prescribed by the
ripple-carry construction, and then appending its two final gates is acyclic: each input
index read by the gate at position $i$ is less than $n+i$.
\end{theorem}

\begin{theorem}[Fan-in of adder gates]
\lean{BoolCircuit.faninTwo_adderGates}
\label{BoolCircuit.faninTwo_adderGates}
\leanok
\uses{BoolCircuit.DAGGate.FaninTwo, BoolCircuit.DAGGate.WellFormed, BoolCircuit.accIdx_lt, BoolCircuit.adderGates, BoolCircuit.faninTwo_adderGatesUpTo}
\difficulty{4}
For every natural number $n$, each gate in the adder circuit with parameter $n$ is well
formed and has at most two inputs.
\end{theorem}

\begin{definition}[Ripple adder circuit]
\lean{BoolCircuit.adderCircuit}
\label{BoolCircuit.adderCircuit}
\leanok
\uses{BoolCircuit.DAGCircuit, BoolCircuit.adderGates, BoolCircuit.gatesAcyclic_adderGates}
For each natural number $n$, define a directed acyclic circuit with the ripple adder
gate list for input size $n$ and output vertex $n+15	ext{⌊}n/3	ext{⌋}+3$.
\end{definition}

\begin{theorem}[Adder circuit has fan-in two]
\lean{BoolCircuit.adderCircuit_isFaninTwo}
\label{BoolCircuit.adderCircuit_isFaninTwo}
\leanok
\uses{BoolCircuit.DAGCircuit.IsFaninTwo, BoolCircuit.adderCircuit, BoolCircuit.faninTwo_adderGates}
\difficulty{1}
For every natural number $n$, the adder circuit on $n$ input bits is well formed, and
each gate reads at most two vertices.
\end{theorem}

\begin{theorem}[Adder circuit size]
\lean{BoolCircuit.adderCircuit_size}
\label{BoolCircuit.adderCircuit_size}
\leanok
\uses{BoolCircuit.DAGCircuit.size, BoolCircuit.adderCircuit, BoolCircuit.adderGates}
\difficulty{3}
For every natural number $n$, the size of the adder circuit on $n$ input vertices is $n
+ 15\lfloor n/3\rfloor + 4$.
\end{theorem}

\begin{theorem}[Adder circuit size at multiples of three]
\lean{BoolCircuit.adderCircuit_size_three_mul}
\label{BoolCircuit.adderCircuit_size_three_mul}
\leanok
\uses{BoolCircuit.DAGCircuit.size, BoolCircuit.adderCircuit, BoolCircuit.adderCircuit_size}
\difficulty{2}
For every natural number $k$, the adder circuit on $3k$ input vertices has size $18k+4$.
\end{theorem}

\begin{definition}[Carry bit]
\lean{BoolCircuit.carryBit}
\label{BoolCircuit.carryBit}
\leanok
Given Boolean sequences $A$ and $B$, define their carry bit at position $0$ to be false;
at position $i+1$, it is true exactly when either both $A_i$ and $B_i$ are true, or
exactly one of them is true and the carry bit at position $i$ is true.
\end{definition}

\begin{definition}[Sum bit]
\lean{BoolCircuit.sumBit}
\label{BoolCircuit.sumBit}
\leanok
\uses{BoolCircuit.carryBit}
Given Boolean sequences $A$ and $B$ and a natural number $i$, define the sum bit at
position $i$ to be the exclusive OR of $A_i$, $B_i$, and their carry bit at position
$i$.
\end{definition}

\begin{definition}[Agreement through a position]
\lean{BoolCircuit.agreeBit}
\label{BoolCircuit.agreeBit}
\leanok
\uses{BoolCircuit.sumBit}
Given Boolean sequences $A$, $B$, and $C$, define agreement through position $0$ to be
true. At position $i+1$, agreement holds exactly when it holds through position $i$ and
the sum bit of $A$ and $B$ at position $i$ equals $C_i$.
\end{definition}

\begin{theorem}[Addition of binary prefix values]
\lean{BoolCircuit.prefixVal_add}
\label{BoolCircuit.prefixVal_add}
\leanok
\uses{BoolCircuit.carryBit, BoolCircuit.prefixVal, BoolCircuit.prefixVal_succ, BoolCircuit.sumBit}
\difficulty{5}
Let $A,B:\mathbb{N}\to\{\mathrm{false},\mathrm{true}\}$ be Boolean sequences, and let
$i$ be a natural number. The sum of their values on positions below $i$, with position
$j$ weighted by $2^j$, equals the corresponding weighted value of their bitwise sum
below $i$, plus $2^i$ if the carry into position $i$ is true (and plus zero otherwise).
\end{theorem}

\begin{theorem}[Agreement on an initial bit segment]
\lean{BoolCircuit.agreeBit_iff}
\label{BoolCircuit.agreeBit_iff}
\leanok
\uses{BoolCircuit.agreeBit, BoolCircuit.sumBit}
\difficulty{5}
For Boolean sequences $A,B,C$ and a natural number $i$, the agreement indicator at $i$
equals true if and only if $C_j$ equals the $j$th sum bit of $A$ and $B$ for every
$j<i$.
\end{theorem}

\begin{theorem}[Prefix addition and absence of carry]
\lean{BoolCircuit.agree_and_noCarry_iff}
\label{BoolCircuit.agree_and_noCarry_iff}
\leanok
\uses{BoolCircuit.agreeBit, BoolCircuit.agreeBit_iff, BoolCircuit.carryBit, BoolCircuit.prefixVal, BoolCircuit.prefixVal_add, BoolCircuit.prefixVal_congr, BoolCircuit.prefixVal_inj, BoolCircuit.prefixVal_lt}
\difficulty{5}
For Boolean sequences $A,B,C$ and a natural number $i$, the conjunction of the agreement
indicator at $i$ with the negation of the carry bit at $i$ equals true if and only if
the value of the first $i$ bits of $C$ is the sum of the values of the first $i$ bits of
$A$ and $B$.
\end{theorem}

\begin{definition}[Input bits with zero padding]
\lean{BoolCircuit.RippleCarry.inBits}
\label{BoolCircuit.RippleCarry.inBits}
\leanok
For $n\in\mathbb{N}$, a Boolean vector
$x\colon\{0,\ldots,n-1\}\to\{\mathrm{false},\mathrm{true}\}$, and $o,j\in\mathbb{N}$,
let $x$ be listed in index order and extended by false values beyond its end. The value
at position $o+j$ in this extended list is the bit $x_{o+j}$ when $o+j<n$, and false
otherwise.
\end{definition}

\begin{theorem}[Input-bit extraction from a list]
\lean{BoolCircuit.RippleCarry.inBits_get}
\label{BoolCircuit.RippleCarry.inBits_get}
\leanok
\uses{BoolCircuit.RippleCarry.inBits}
\difficulty{1}
For every Boolean list $w$ and natural numbers $o,j$, extracting the input bit at offset
$o+j$ from the sequence of entries of $w$ gives $w[o+j]$ when that position exists, and
false otherwise.
\end{theorem}

\begin{theorem}[Ripple-carry circuit specification]
\lean{BoolCircuit.RippleCarry.adderGatesUpTo_spec}
\label{BoolCircuit.RippleCarry.adderGatesUpTo_spec}
\leanok
\uses{BoolCircuit.DAGGate.eval, BoolCircuit.RippleCarry.accIdx, BoolCircuit.RippleCarry.adderGatesUpTo, BoolCircuit.RippleCarry.cinIdx, BoolCircuit.RippleCarry.inBits, BoolCircuit.RippleCarry.runWith_bitBlock, BoolCircuit.accIdx_lt, BoolCircuit.agreeBit, BoolCircuit.carryBit, BoolCircuit.cinIdx_lt, BoolCircuit.constGate, BoolCircuit.runWith, BoolCircuit.runWith_append, BoolCircuit.runWith_cons, BoolCircuit.runWith_getD_of_lt}
\difficulty{6}
Let $n$ be a natural number, let $x$ assign a Boolean value to each of the $n$ input
positions, and let $k,i$ be natural numbers with $3k\le n$ and $i\le k$. After running
the circuit consisting of the initial two constant gates followed by $i$ bit blocks on
the input list of $x$, the resulting list has length $n+2+15i$. Its value at the carry
index is the carry bit after $i$ positions for the input streams beginning at offsets
$0$ and $k$. Its value at the accumulator index is true exactly when, at every position
below $i$, the sum bit from the streams beginning at offsets $0$ and $k$ agrees with the
input stream beginning at offset $2k$.
\end{theorem}

\begin{theorem}[Adder circuit evaluation]
\lean{BoolCircuit.adderCircuit_eval}
\label{BoolCircuit.adderCircuit_eval}
\leanok
\uses{BoolCircuit.DAGCircuit.eval, BoolCircuit.DAGCircuit.values, BoolCircuit.DAGGate.eval, BoolCircuit.RippleCarry.adderGatesUpTo, BoolCircuit.RippleCarry.adderGatesUpTo_spec, BoolCircuit.RippleCarry.inBits, BoolCircuit.accIdx_lt, BoolCircuit.adderCircuit, BoolCircuit.adderGates, BoolCircuit.agreeBit, BoolCircuit.carryBit, BoolCircuit.runWith, BoolCircuit.runWith_append, BoolCircuit.runWith_cons, BoolCircuit.runWith_nil}
\difficulty{4}
For every natural number $n$ and every Boolean assignment $x$ to $n$ input bits, the
adder circuit's output equals the conjunction of the following two conditions: the first
$n/3$ bits of $x$, the next $n/3$ bits, and the subsequent $n/3$ bits agree bit by bit
with the sum bits produced by binary addition of the first two blocks; and the final
carry from adding the first two blocks is false.
\end{theorem}

\begin{theorem}[Adder circuit recognizes binary addition]
\lean{BoolCircuit.adderCircuit_eval_append}
\label{BoolCircuit.adderCircuit_eval_append}
\leanok
\uses{BoolCircuit.DAGCircuit.eval, BoolCircuit.RippleCarry.inBits, BoolCircuit.RippleCarry.inBits_get, BoolCircuit.adderCircuit, BoolCircuit.adderCircuit_eval, BoolCircuit.agree_and_noCarry_iff, BoolCircuit.bitsVal, BoolCircuit.bitsVal_eq_prefixVal, BoolCircuit.prefixVal_congr}
\difficulty{5}
Let $a$, $b$, and $c$ be finite Boolean lists of equal length. The adder circuit on the
concatenated input $a,b,c$ outputs true if and only if the integer represented by $c$ in
binary equals the sum of the integers represented by $a$ and $b$ in binary.
\end{theorem}
